% from ega4-11-fr.tex
\oldpage[IV-3]{153}
fermé,
l'homomorphisme $\mathcal O_z\to(h_*\mathcal O_{Z'})_z$ est injectif, car il
résulte de (2.3.1) que $h_*\mathcal O_{Z'}=g_*\mathcal O_{Y'}\otimes_Y Z$,
et notre assertion résulte de l'hypothèse (i) et de la définition des modules
plats $(\mathbf 0_{\mathrm I},6.1.1)$.

On peut donc désormais supposer l'anneau local noethérien $A$ complet~; la
démonstration sera achevée si l'on prouve que l'intersection des noyaux des
homomorphismes $A=\mathcal O_y\to\mathcal O_{y'}$ (où $y'$ parcourt
$g^{-1}(y)$) est réduite à $0$. En effet, les $\mathcal O_{y'}$ sont des anneaux
locaux noethériens, donc pour chaque $y'$ l'intersection des
$\mathfrak m_{y'}^n$ $(n\geq0)$ est réduite à $0$~; si $\mathfrak a_{n,y'}$ est
l'image réciproque dans $A$ de $\mathfrak m_{y'}^n$, les intersections finies
des $\mathfrak a_{n,y'}$ sont des voisinages de $0$ dans $A$ et l'intersection de
tous les $\mathfrak a_{n,y'}$ est réduite à $0$~; elles forment donc un système
fondamental de voisinages de $0$ dans $A$ (Bourbaki, \emph{Alg. comm.}, chap.~III,
\S~2, n$^\circ$~7, prop.~8, où l'on peut remplacer la suite décroissante par un
ensemble filtrant quelconque), et l'on pourra par suite appliquer (11.5.1).
Or, soit $s\in A$ un élément appartenant au noyau de chacun des homomorphismes
$A\to\mathcal O_{y'}$~; l'image $s'$ de $s$ dans
$\Gamma(Y',\mathcal O_{Y'})$ est donc une section de $\mathcal O_{Y'}$ au-dessus
de $Y'$ telle que $s'_{y'}=0$ pour tout $y'\in g^{-1}(y)$~; il y a par suite un
voisinage $V'$ de $g^{-1}(y)$ dans $Y'$ tel que $s'|V'=0$. Mais comme $g$ est
fermé, $V'$ contient un ensemble de la forme $g^{-1}(V)$, où $V$ est un
voisinage ouvert de $y$ dans $Y$~; or $Y$ est un schéma local, donc le seul
voisinage du point fermé $y$ de $Y$ est $Y$ tout entier, autrement dit $V'=Y'$,
$s'=0$, et comme $A\to\Gamma(Y',\mathcal O_{Y'})$ est injectif par hypothèse,
$s=0$. C.Q.F.D.
\end{proof}

Le cas particulier suivant de (11.5.3) nous sera utile au chap.~V.

\begin{corollary}[11.5.4]
\label{IV.11.5.4-fr}
Soit $f=(\psi,\theta):X\to Y$ un morphisme propre de présentation finie, et
soit $p:X\times_YX\to X$ la première projection. On suppose que
$\theta:\mathcal O_Y\to f_*(\mathcal O_X)$ est injectif. Alors, pour que $f$
soit plat, il faut et il suffit que $p$ le soit.
\end{corollary}

\begin{proposition}[11.5.5]
\label{IV.11.5.5-fr}
Soient $A$ un anneau, $Y=\operatorname{Spec}(A)$, $f:X\to Y$ un morphisme
localement de présentation finie, $\mathcal F$ un $\mathcal O_X$-Module
quasi-cohérent de présentation finie, $x$ un point de $X$. Soit
$A\to A'$ un homomorphisme injectif faisant de $A'$ une algèbre entière sur $A$.
Posons $Y'=\operatorname{Spec}(A')$, $X'=X\times_Y Y'$, $f'=f\times1$,
$\mathcal F'=\mathcal F\otimes_Y Y'$. Alors, si $\mathcal F'$ est $f'$-plat en
tout point de $X'$ dont la projection dans $X$ est égale à $x$, $\mathcal F$ est
$f$-plat au point $x$.
\end{proposition}

\begin{proof}
Supposons d'abord que $A'$ soit une $A$-algèbre finie et de présentation finie~;
alors le morphisme $Y'\to Y$ est propre (\textbf{II}, 6.1.11) et de présentation
finie, donc les hypothèses de (11.5.3) sont vérifiées, d'où la conclusion. Dans
le cas général, la proposition résultera de ce cas particulier, du fait que $A'$
est limite inductive de ses sous-$A$-algèbres finies $A'_\lambda$ et des deux
lemmes suivants~:
\end{proof}

\begin{lemma}[11.5.5.1]
\label{IV.11.5.5.1-fr}
Toute $A$-algèbre finie $A'$ est limite inductive de $A$-algèbres $A'_\lambda$
qui sont finies et de présentation finie.
\end{lemma}

\begin{proof}
On raisonne comme dans (1.9.3.1). On a en effet $A'=B/\mathfrak J$, où $B$ est
une $A$-algèbre finie qui est un $A$-module libre, et $\mathfrak J$ un idéal de
$B$ (1.4.7.1). Or, $\mathfrak J$ est limite inductive des idéaux
$\mathfrak J_\lambda\subset\mathfrak J$ de $B$ qui sont de type fini (et a
fortiori des $A$-modules de type fini), donc, par l'exactitude du foncteur
$\varinjlim$, $A'$ est limite inductive des $A$-algèbres
$B/\mathfrak J_\lambda$~; or $B/\mathfrak J_\lambda$ est par définition un
$A$-module de présentation finie, donc aussi (1.4.7) une $A$-algèbre de
présentation finie.
\end{proof}
% IV3_B22_P153_END

